\documentclass[12pt]{article}
\usepackage{geometry}
\geometry{margin=1in}
\usepackage{setspace}
\usepackage{array}
\usepackage{booktabs}

\setstretch{1.2}

\title{Quiz 3\\ \vspace{8pt} \Large ECON 308 -- International Economic Relations}
\author{Instructor: Muhebullah Karimzada}
\date{October 21,  2025}
\usepackage{comment}

\begin{document}
\maketitle


\section*{Multiple Choice Questions (1 point each)}

\begin{enumerate}
    \item According to the Heckscher--Ohlin model, differences in trade patterns across countries arise primarily due to differences in:  
    \begin{enumerate}
        \item Technology  
        \item Preferences  
        \item Factor endowments  
        \item Exchange rates  
    \end{enumerate}

    \item Suppose Canada is capital-abundant and Mexico is labor-abundant. According to the Heckscher--Ohlin theorem, after trade opens:  
    \begin{enumerate}
        \item Canada exports labor-intensive goods  
        \item Canada exports capital-intensive goods  
        \item Mexico exports capital-intensive goods  
        \item Both countries export the same goods  
    \end{enumerate}

    \item The Stolper--Samuelson theorem states that:  
    \begin{enumerate}
        \item Trade benefits all factors of production equally  
        \item Trade hurts both factors in the import-competing sector  
        \item Trade raises the real income of the abundant factor and lowers that of the scarce factor  
        \item Trade raises the real income of labor but reduces that of capital  
    \end{enumerate}

    \item The Rybczynski theorem implies that, holding goods prices constant, an increase in a country's endowment of labor will:  
    \begin{enumerate}
        \item Increase output of both goods  
        \item Increase output of the labor-intensive good and decrease output of the capital-intensive good  
        \item Decrease output of both goods  
        \item Have no effect on production  
    \end{enumerate}

    \item The Leontief Paradox showed that:  
    \begin{enumerate}
        \item The U.S. exported capital-intensive goods as predicted  
        \item The U.S. imported labor-intensive goods despite being capital-abundant  
        \item The Heckscher--Ohlin model perfectly predicted trade patterns  
        \item The U.S. was actually labor-abundant in 1947  
    \end{enumerate}
\end{enumerate}

\vspace{0.5cm}
\newpage
\section*{Problem (15 points)}

\noindent Consider two countries, \textbf{Home (H)} and \textbf{Foreign (F)}, that both produce two goods: \textbf{Computers (C)} (capital-intensive) and \textbf{Textiles (T)} (labor-intensive). Both countries have identical technologies and homothetic preferences.

\vspace{0.4cm}

\noindent \textbf{Factor Endowments:}
\begin{itemize}
    \item Home: $K_H = 300$ units of capital, $L_H = 300$ units of labor
    \item Foreign: $K_F = 100$ units of capital, $L_F = 300$ units of labor
\end{itemize}

\vspace{0.3cm}

\noindent \textbf{Production Requirements (per unit of output):}

\begin{center}
\begin{tabular}{lcc}
\toprule
\textbf{Good} & \textbf{Capital per unit ($a_K$)} & \textbf{Labor per unit ($a_L$)} \\
\midrule
Computers (C) & 6 & 3 \\
Textiles (T) & 2 & 4 \\
\bottomrule
\end{tabular}
\end{center}

\vspace{0.5cm}


\begin{enumerate}
    \item[(a)] (3 pts) Determine which country is capital-abundant and which is labor-abundant.  
    \item[(b)] (3 pts) Identify which good is capital-intensive.  
    \item[(c)] (4 pts) Using the Heckscher--Ohlin theorem, predict the trade pattern between Home and Foreign.  
    \item[(d)] (5 pts) Using the Stolper--Samuelson theorem, explain what happens to the real returns to capital and labor in Home after trade opens.  
\end{enumerate}

\vfill
\noindent \textbf{Total: 20 Points}

\end{document}
